\section{Benchmark Implications}
% ══════════════════════════════════════════════════════════════

The results suggest three rules for future cross-sectional LLM ranker benchmarks.
First, baselines should include factor sorts and lightweight supervised rankers.
Factor sorts test whether an LLM is rediscovering familiar style premia, while expanding-window models test whether the same panel already contains enough signal for a conventional learner.
The gap between Random Forest/Ridge and the strongest LLMs in Tables~\ref{tab:longshort}--\ref{tab:ff6} is therefore more informative than a comparison against the market index alone.

Second, evaluation should keep the ranking object fixed.
Pointwise scores, pairwise tournaments, setwise reranking, and agentic debate can all be useful, but they should ultimately output the same monthly ordering over the same investable universe.
This makes IC, quintile monotonicity, long--short spreads, and factor diagnostics directly comparable across methods.
It also prevents an agentic system from changing the task by using additional data, extra lookups, or different rebalancing assumptions.

Third, cost and reproducibility should be first-class metrics.
Listwise prompting is attractive because it obtains the full cross-section in one call, whereas pairwise and agentic systems can scale call counts quickly.
Reporting cost per monthly experiment, malformed-output handling, decoding settings, and conservative evaluation boundaries makes it possible to distinguish genuine ranking skill from engineering choices that are expensive or difficult to replicate.

These rules also clarify how to read the negative findings.
The underperforming models are not simply lower-return versions of the winners; they encode different implicit priors about the cross-section.
GPT-4o mini, for example, ranks toward reversal and high-volatility names in a period when the opposite style was rewarded.
A future benchmark should therefore measure not only whether a model wins in one sample but also whether its style exposure is stable, interpretable, and robust when market leadership changes.

Finally, the current evidence should be viewed as a controlled starting point rather than a finished trading system.
The 30-stock universe is small enough for transparent listwise prompting, and the transaction-cost model is intentionally simple.
Broader universes, rolling sector composition, liquidity constraints, and turnover-aware construction are natural next steps.
Future work should add broader universes, rolling membership, and richer information axes---such as disclosure text, which recent real-time event benchmarks show can explain announcement returns \citep{koijen2026agentic}---while keeping raw ranking quality separate from weighting or optimization gains.

% ══════════════════════════════════════════════════════════════
